\section{Reproducibility Details}
\paragraph{Execution context.}
The second-pass results used in the main paper were produced locally under
Python 3.11.14 on macOS arm64 (Apple Silicon, 8 logical CPU cores) with
PyTorch 2.11.0, Gymnasium 1.2.3, and Minigrid 3.0.0. The preflight audit
selected the repository's \texttt{PROFILE\_FULL\_SAFE} mode and recorded the
exact package versions before the second-pass rerun.

\paragraph{Exact commands.}
The submission supplement includes an anonymized executable artifact containing
the benchmark code, saved outputs, and paper assets. The main reproduction path
is:
\begin{itemize}
\item \texttt{make preflight}
\item \texttt{make stage1}
\item \texttt{./.venv/bin/python scripts/run\_second\_pass.py}
\item \texttt{./.venv/bin/python scripts/make\_plots.py --results-root results/second\_pass}
\item \texttt{./.venv/bin/python scripts/make\_tables.py --results-root results/second\_pass}
\item \texttt{./.venv/bin/python scripts/make\_paper.py --results-root results/second\_pass --latex-root neurips\_latex}
\end{itemize}
These commands rebuild the figures, tables, and populated result sections from
the saved artifact roots used in the paper.

\paragraph{Seed variability and ranking stability.}
Table~\ref{tab:seed-variability} reports ID seed variability for the main proxy
and oracle quantities. We do not claim tight statistical confidence from only
three seeds, but the benchmark ranking is still stable in a stronger sense:
for every environment and both ID/OOD splits,
\[
\max_{\text{seed}} \; \text{true\_score}(\text{Proxy PPO})
<
\min_{\text{seed}} \; \text{true\_score}(\text{Oracle True PPO}).
\]
Thus the proxy-versus-oracle gap is not a single-seed artifact.

\begin{table}[h]
\centering
\begin{tabular}{l c c c}
\toprule
Environment & Proxy ID true & Proxy ID hack & Oracle ID true \\
\midrule
ProxyFarming & -1.950 +/- 0.000 & 1.000 +/- 0.000 & 0.992 +/- 0.014 \\
SensorTamper & -3.346 +/- 1.285 & 0.833 +/- 0.289 & 1.000 +/- 0.000 \\
GemTimeout & -0.153 +/- 0.053 & 1.000 +/- 0.000 & 1.000 +/- 0.000 \\
\bottomrule
\end{tabular}
\caption{ID seed variability across the three proxy and oracle runs. Entries are mean +/- standard deviation over seeds.}
\label{tab:seed-variability}
\end{table}


\paragraph{Compute and storage.}
The full second-pass rerun completed in 102.2 seconds wall-clock. The only
retrained environment in that pass was GemTimeout-v0, which consumed 60,000 PPO
steps per seed and 92.5 seconds total wall-clock across Proxy PPO and Oracle
True PPO. ProxyFarming-v0 and SensorTamper-v0 reused existing checkpoints and
were re-evaluated only. The second-pass GemTimeout retraining artifacts occupy
roughly 3.0 MB across the two methods.

\paragraph{Assets and licenses.}
The executable benchmark reuses upstream open-source software whose versions and
licenses are recorded in the local environment: PyTorch 2.11.0
(BSD-3-Clause), Gymnasium 1.2.3 (MIT), and Minigrid 3.0.0 (MIT). The new assets
introduced by this work are the benchmark environments, schema exports, saved
result bundle, plotting and table-generation scripts, and the NeurIPS paper
package. Those assets are documented by the repository \texttt{README.md},
\texttt{neurips\_latex/BUILD.md}, schema definitions, and the generated result
inventory under \texttt{results/}.
